\[
\vect u=\text{constante}
\]
cambia la posición de todos los puntos pero no las distancias relativas entre ellos.

También una pequeña rotación rígida puede cambiar la orientación sin producir deformación.

\begin{idea}
La deformación mide el \textbf{cambio relativo entre puntos próximos}, no el movimiento absoluto del cuerpo.
\end{idea}

\safesection{Deformación longitudinal}
Para una barra de longitud inicial $L_0$ que pasa a medir $L$:
\[
\Delta L=L-L_0.
\]

La deformación longitudinal nominal es
\[
\varepsilon=\frac{\Delta L}{L_0}
\]
y es adimensional.

Si $\varepsilon>0$ hay alargamiento. Si $\varepsilon<0$ hay acortamiento.

\safesection{Deformación verdadera}
La deformación verdadera incremental se define como
\[
\dd\varepsilon_{\mathrm{real}}=\frac{\dd L}{L}.
\]
Integrando:
\[
\varepsilon_{\mathrm{real}}=\ln\left(\frac{L}{L_0}\right).
\]
Para deformaciones pequeñas:
\[
\ln(1+x)\simeq x,
\]
y, por tanto,
\[
\varepsilon_{\mathrm{real}}\simeq\varepsilon_{\mathrm{nom}}.
\]

\safesection{Tensor de pequeñas deformaciones}
Sea el campo de desplazamientos
\[
\vect u=(u_1,u_2,u_3).
\]
En pequeñas deformaciones:
\[
\varepsilon_{ij}
=\frac{1}{2}
\left(
\frac{\partial u_i}{\partial x_j}
+
\frac{\partial u_j}{\partial x_i}
\right).
\]

De forma compacta:
\[
\mat{\varepsilon}
=\frac12
\left[
\nabla\vect u+(\nabla\vect u)^T
\right].
\]

En dos dimensiones:
\[
\varepsilon_{xx}=\frac{\partial u}{\partial x},\qquad
\varepsilon_{yy}=\frac{\partial v}{\partial y},
\]
\[
\varepsilon_{xy}
=\frac12\left(
\frac{\partial u}{\partial y}
+
\frac{\partial v}{\partial x}
\right).
\]

La deformación angular ingenieril es
\[
\gamma_{xy}=2\varepsilon_{xy}.
\]

\begin{errorbox}
No confundas $\gamma_{xy}$ con $\varepsilon_{xy}$. En la convención tensorial:
\[
\gamma_{xy}=2\varepsilon_{xy}.
\]
\end{errorbox}

% ============================================================
\safechapter{Leyes de comportamiento}
% ============================================================

\safesection{Ley constitutiva}
El equilibrio y la cinemática no bastan.

Necesitamos relacionar
\[
\text{tensiones}\longleftrightarrow\text{deformaciones}
\]
mediante una ley constitutiva.

En elasticidad lineal general:
\[
\sigma_{ij}=C_{ijkl}\varepsilon_{kl},
\]
donde $C_{ijkl}$ contiene las propiedades elásticas del material.

\safesection{Ley de Hooke uniaxial}
Para una barra sometida a tensión uniaxial dentro del régimen elástico lineal:
\[
\sigma=E\varepsilon,
\]
donde $E$ es el módulo de Young.

